% STATUS: written 2026-11-05 (Cheung 2001 pp. 117--127 read in full; Cheung 2003
% pp. 299--310 read in full; Benkovic 2014 Section 1 and Theorem 1.1; Wang 2016
% pp. 1952--1953 read for the triangular-ring form, already used in chapter 7.)
% STATUS: extended 2026-11-06 (Vladeva 2018 read in full)
% STATUS: extended 2026-11-07 (Cheng and Jing 2007, Ma and Ji 2007, Xiao 2012,
% Benkovic 2011, Liang--Ren--Wei 2020 and Cheraghpour--Jafari 2024 read in full)
% Sources: Vladeva 2018 (pp. 1--5); Cheng and Jing 2007 (pp. 1--5);
% Ma and Ji 2007 (pp. 355--362); Xiao 2012 (pp. 1--7); Benkovic 2011
% (pp. 1532--1541); Liang, Ren and Wei 2020 (pp. 1--8); Cheraghpour and
% Jafari 2024 (pp. 1--6).
\section{Triangular algebras}\label{sec:matrix-triangular:triangular}

Section~\ref{sec:matrix-triangular:matrix} treated full matrix algebras, where the
two algebras $A$ and $B$ of the block picture have the same size and the bimodule
$M$ fills the whole off-diagonal block \cite[p.~119]{Cheung2001}. A triangular
algebra keeps the block shape and sets one off-diagonal block to zero
\cite[p.~119]{Cheung2001}. That is the smallest change which removes the symmetry
that the matrix-unit arguments used \cite[p.~119]{Cheung2001}.

\begin{definition}[{\cite[Definition~2, p.~119]{Cheung2001}}]\label{def:triangular:algebra}
Let $A$ and $B$ be unital algebras over a commutative ring $R$, and let $M$ be an
$(A,B)$-bimodule that is faithful as a left $A$-module and as a right
$B$-module. The $R$-algebra
\[
\operatorname{Tri}(A,M,B)=\left\{\begin{pmatrix} a & m\\ 0 & b\end{pmatrix}:
a\in A,\ m\in M,\ b\in B\right\}
\]
under the usual matrix addition and formal matrix multiplication is called a
triangular algebra. We write $a\oplus b$ for the element with $m=0$.
\end{definition}

Cheung notes that the two natural projections $\pi_A$ and $\pi_B$, which read off
$a$ and $b$, are the only maps needed to state the results
\cite[Definition~3, p.~119]{Cheung2001}. Cheung gives the centre a short description, and
that description carries the whole theory (Proposition~\ref{prop:triangular:centre}) \cite[p.~119]{Cheung2001}.

\begin{proposition}[{\cite[Proposition~3, p.~119]{Cheung2001}}]\label{prop:triangular:centre}
The centre of $\Omega=\operatorname{Tri}(A,M,B)$ is
\[
Z(\Omega)=\{a\oplus b : a\in A,\ b\in B,\ am=mb\ \text{for every } m\in M\}.
\]
Moreover $\pi_A(Z(\Omega))\subseteq Z(A)$ and $\pi_B(Z(\Omega))\subseteq Z(B)$,
and there is a unique algebra isomorphism $\tau$ from $\pi_A(Z(\Omega))$ to
$\pi_B(Z(\Omega))$ such that $am=m\tau(a)$ for every $m\in M$.
\end{proposition}

The faithfulness of $M$ on both sides enters here. A central element $a\oplus b$
forces $a\in Z(A)$ and $b\in Z(B)$, since $M$ separates $A$ from $B$
\cite[p.~119]{Cheung2001}.

Derivations of a triangular algebra split in an explicit way. Cheung records the
split as a lemma, and the rest of the paper uses it (Lemma~\ref{lem:triangular:derivation})
\cite[Lemma~5, p.~303]{Cheung2003}.\begin{lemma}[{\cite[Lemma~5, p.~303]{Cheung2003}}]\label{lem:triangular:derivation}
A derivation $\delta$ on $\Omega=\operatorname{Tri}(A,M,B)$ is of the form
\[
\delta\begin{pmatrix} a & m\\ 0 & b\end{pmatrix}
=\begin{pmatrix} p_A(a) & an-nb+f(m)\\ 0 & p_B(b)\end{pmatrix},
\]
where $n\in M$ and
\begin{enumerate}
\item $p_A$ is a derivation on $A$ and $f(am)=p_A(a)m+af(m)$;
\item $p_B$ is a derivation of $B$ and $f(mb)=mp_B(b)+f(m)b$.
\end{enumerate}
\end{lemma}

The element $n$ gives the inner part of the derivation. The map $f$ is the part
that acts on the bimodule \cite[p.~303]{Cheung2003}. Suppose $M$ is faithful on
the left. Then the first condition alone forces $p_A$ to be a derivation, so the
two conditions are not independent \cite[p.~303]{Cheung2003}.

One more name is needed, for the maps that the derivation lemma does not cover
\cite[p.~299]{Cheung2003}.

\begin{definition}[{\cite[p.~299]{Cheung2003}}]\label{def:triangular:proper-lie}
A proper Lie derivation is a Lie derivation of an associative algebra that can be
written as the sum of a derivation and a linear map with image in the centre of the
algebra. A Lie derivation that is not a proper Lie derivation is called an improper
Lie derivation.
\end{definition}

Cheung first describes the shape of an arbitrary Lie derivation
\cite[Proposition~4, p.~302]{Cheung2003}. The description has one element
$n\in M$ and six linear maps \cite[Proposition~4, p.~302]{Cheung2003}. A map is
proper exactly when two centre-valued maps $h_A$ and $h_B$ pass a test
\cite[Theorem~6, p.~303]{Cheung2003}.

The main theorem asks when every Lie derivation is proper. Its conditions
involve only $A$ and $B$ \cite[Theorem~11, p.~307]{Cheung2003}. For an algebra
$X$, let $W(X)$ be the smallest subalgebra that meets conditions (i) to (vii) of
Proposition~10 of that paper \cite[Proposition~10, p.~305]{Cheung2003}.

\begin{theorem}[{\cite[Theorem~11, p.~307]{Cheung2003}}]\label{thm:triangular:proper-lie}
Every Lie derivation of $\Omega=\operatorname{Tri}(A,M,B)$ is proper if the
following two conditions hold:
\begin{enumerate}
\item either $M$ is a faithful left $A$-module and $Z(B)=\pi_B(Z(\Omega))$; or $M$
is a faithful left $A$-module and $A=W(A)$; or every Lie derivation of $A$ is
proper and $A=W(A)$;
\item either $M$ is a faithful right $B$-module and $Z(A)=\pi_A(Z(\Omega))$; or $M$
is a faithful right $B$-module and $B=W(B)$; or every Lie derivation of $B$ is
proper and $B=W(B)$.
\end{enumerate}
\end{theorem}

Cheung gives three simple ways to meet $A=W(A)$. One is that $A$ is generated
by commutators and idempotents. Another is that two elements of $A$ commute to
an invertible element \cite[p.~307]{Cheung2003}. The upper triangular matrix
algebra meets the first one, which gives the next corollary (Corollary~\ref{cor:triangular:upper-lie}).

\begin{corollary}[{\cite[Corollary~17, p.~308]{Cheung2003}}]\label{cor:triangular:upper-lie}
Every Lie derivation of $T_n(R)$, the $n\times n$ upper triangular matrix algebra
over $R$, has the form $a\mapsto ax-xa+h(a)1$, where $x\in T_n(R)$ is fixed and
$h$ is a linear functional defined on the span of $\{E_{jj}:j=1,\dots,n\}$.
\end{corollary}

Over a semiring the upper triangular matrix algebra still has derivations
\cite[p.~5]{Vladeva2018Semirings}. Take an additively idempotent semiring, where
the sum of two equal terms adds nothing \cite[p.~2]{Vladeva2018Semirings}. Set
one principal dense submatrix to zero and keep every other entry. That map is a
derivation \cite[Theorem~1, p.~5]{Vladeva2018Semirings}. The partial identity
matrices produce these maps, and the theorem counts them
\cite[p.~5]{Vladeva2018Semirings}.

The commuting-map question has a parallel answer on the same class of algebras
\cite[Theorem~2, p.~124]{Cheung2001}. Cheung proves that every commuting map is
proper under three conditions \cite[Theorem~2, p.~124]{Cheung2001}. Each centre
must be the image of the centre of $\Omega$. One element $m_0\in M$ must
separate the central elements \cite[Theorem~2, p.~124]{Cheung2001}. The proof
first describes an arbitrary commuting map. Its data are two commuting maps, one
on $A$ and one on $B$, and four maps into centres
\cite[Proposition~4, p.~121]{Cheung2001}.

The conditions are needed. A triangular algebra can carry an improper commuting
map even when $A$ and $B$ are both commutative
\cite[Example~1, p.~123]{Cheung2001}. That example is why the chapter keeps the
commuting-map and the derivation theorems apart
\cite[Example~1, p.~123]{Cheung2001}.

The Jordan law holds on triangular algebras without the unital hypothesis. One
condition on the two diagonal algebras is enough
\cite[Theorem~2.10, p.~5]{ChengJing2007}. Cheng and Jing say that $A$ has
property (P) when $xay+yax=0$ for all $x,y\in A$ forces $a=0$ (Theorem~\ref{thm:triangular:cheng-jing})
\cite[p.~2]{ChengJing2007}.\begin{theorem}[{\cite[Theorem~2.10, p.~5]{ChengJing2007}}]\label{thm:triangular:cheng-jing}
Let $A$ and $B$ be two algebras over a $2$-torsion free commutative ring $R$,
both with property (P), and let $M$ be a faithful $(A,B)$-bimodule. Then every
Jordan derivation of $\operatorname{Tri}(A,M,B)$ into itself is a derivation.
\end{theorem}

Neither $A$ nor $B$ need have an identity, and the proof does not use the two
idempotents $P$ and $Q$ of the block picture \cite[p.~1]{ChengJing2007}. The
authors prove the three component statements separately: $\delta$ is a
derivation on each of $T_{11}$, $T_{12}$ and $T_{22}$
\cite[Lemmas~2.7--2.9, pp.~4--5]{ChengJing2007}.

The same law reaches a bimodule when the map is generalized by its value at the
identity \cite[Theorem~2.4, p.~362]{MaJi2007}. A generalized Jordan derivation
$\Phi$ satisfies
$\Phi(ab+ba)=\Phi(a)b+a\Phi(b)+\Phi(b)a+b\Phi(a)-a\Phi(I)b-b\Phi(I)a$
\cite[p.~355]{MaJi2007}. A generalized derivation $\Psi$ satisfies
$\Psi(ab)=\Psi(a)b+a\Psi(b)-a\Psi(I)b$ (Theorem~\ref{thm:triangular:ma-ji}) \cite[p.~355]{MaJi2007}.

\begin{theorem}[{\cite[Theorem~2.4, p.~362]{MaJi2007}}]\label{thm:triangular:ma-ji}
Let $\Phi$ be a generalized Jordan derivation from $T_n(R)$ into a
$T_n(R)$-bimodule $M$. Then there are a unique generalized derivation $\Psi$ and
an antiderivation $\Lambda$ from $T_n(R)$ to $M$, with
$\Lambda(D_n(R))=0$, such that $\Phi=\Psi+\Lambda$.
\end{theorem}

Here $D_n(R)$ is the subalgebra of diagonal matrices
\cite[p.~356]{MaJi2007}. The condition on $D_n(R)$ makes the factorization
unique. Without it an antiderivation can appear twice
\cite[p.~362]{MaJi2007}.

The Jordan question also has a nonlinear form. A nonlinear Jordan derivation is
a map that satisfies the Jordan law with no additivity assumed
\cite[p.~1]{Xiao2012}. Xiao shows that on a triangular algebra the nonlinear law
forces additivity and the derivation rule together (Theorem~\ref{thm:triangular:nonlinear-jordan})
\cite[Theorem~2.1, p.~2]{Xiao2012}.\begin{theorem}[{\cite[Theorem~2.1, p.~2]{Xiao2012}}]\label{thm:triangular:nonlinear-jordan}
Let $T=\begin{pmatrix}A&M\\0&B\end{pmatrix}$ be a triangular algebra and let $d$
be a nonlinear Jordan derivation of $T$. Then $d$ is an additive derivation.
\end{theorem}

The proof first finds the image of each corner. The image of $T_{12}$ lies in
$T_{12}$, and the image of $T_{11}$ lies in $T_{11}+T_{12}$
\cite[Lemmas~2.2--2.3, p.~2]{Xiao2012}. The nonlinear hypothesis then gives
additivity, and the derivation rule follows \cite[pp.~2--7]{Xiao2012}. Take a
nest algebra on an infinite dimensional Hilbert space. There every nonlinear
Jordan derivation is inner. It is an additive derivation, and every derivation
of a nest algebra is inner \cite[Proposition~2.11, p.~7]{Xiao2012}.

Lie derivations of a triangular algebra carry a generalization of their own. A
generalized Lie derivation $f$ comes with an associated linear map $d$ and
satisfies $f([x,y])=[f(x),y]+[x,d(y)]$ \cite[p.~1535]{Benkovic2011GenLie}.
Benkovic reduces the question to a condition on the single element $d(1)$ (Theorem~\ref{thm:triangular:generalized-lie})
\cite[p.~1539]{Benkovic2011GenLie}.\begin{theorem}[{\cite[Theorem~3.6, p.~1541]{Benkovic2011GenLie}}]\label{thm:triangular:generalized-lie}
Let $\mathcal{A}=\operatorname{Tri}(A,M,B)$ be a triangular algebra and let
$f\colon\mathcal{A}\to\mathcal{A}$ be a generalized Lie derivation with
associated linear map $d$. Suppose that
\begin{enumerate}
\item $\pi_A(Z(\mathcal{A}))=Z(A)$ or $[B,B]=B$, and
\item $\pi_B(Z(\mathcal{A}))=Z(B)$ or $[A,A]=A$.
\end{enumerate}
Then $f=\Delta+\tau$, where $\Delta$ is a generalized derivation of
$\mathcal{A}$ and $\tau\colon\mathcal{A}\to Z(\mathcal{A})$ is a linear map
that vanishes on all commutators of $\mathcal{A}$.
\end{theorem}

The two conditions let the theorem of Cheung on proper Lie derivations reach the Lie
part \cite[p.~1541]{Benkovic2011GenLie}. Take a Hilbert space nest algebra. Its
centre is one-dimensional. There every generalized Lie derivation is a
generalized inner derivation plus a centre-valued linear map
\cite[Corollary~3.8, p.~1542]{Benkovic2011GenLie}. The same paper settles the
generalized Jordan question on a $2$-torsion free triangular algebra
\cite[Theorem~4.2, p.~1542]{Benkovic2011GenLie}.

The bilinear form of the question has the same shape. A biderivation of an
algebra $A$ is a bilinear map that is a derivation in each argument
\cite[p.~2]{LiangRenWei2020}. Its Lie form is a Lie derivation in each argument (Theorem~\ref{thm:triangular:lie-biderivation})
\cite[p.~2]{LiangRenWei2020}.\begin{theorem}[{\cite[Theorem~3.2, p.~7]{LiangRenWei2020}}]\label{thm:triangular:lie-biderivation}
Let $T=\begin{pmatrix}A&M\\0&B\end{pmatrix}$ be a triangular algebra over a
commutative ring $R$ and let $\varphi\colon T\times T\to T$ be a Lie
biderivation. Suppose that
\begin{enumerate}
\item $\pi_A(Z(T))=Z(A)$ and $\pi_B(Z(T))=Z(B)$;
\item at least one of $A$ and $B$ is noncommutative;
\item every bimodule homomorphism $f\colon M\to M$ has the standard form;
\item $\alpha a=0$ with $\alpha\in Z(A)$ and $0\neq a\in A$ forces $\alpha=0$.
\end{enumerate}
Then $\varphi(x,y)=\lambda_0[x,y]+[x,[y,\varphi(e,e)]]+\mu(x,y)$ for some
$\lambda_0\in Z(T)$ and some central bilinear map $\mu$.
\end{theorem}

A bimodule homomorphism has the standard form when
$f(m)=a_0m+mb_0$ for central $a_0$ and $b_0$
\cite[equation~(3.3), p.~7]{LiangRenWei2020}. On a block upper triangular matrix
algebra over a domain the sum reduces to one extremal biderivation plus one inner
biderivation \cite[Corollary~3.8, p.~8]{LiangRenWei2020}.

A superalgebra version of the triangular ring has the same block shape with a
sign in the product \cite[p.~1]{CheraghpourJafari2024}. A Jordan superderivation
of degree $0$ satisfies the Jordan law, and one of degree $1$ satisfies it with
the super sign \cite[p.~2]{CheraghpourJafari2024}. The upper triangular matrix
ring here is the trivial extension of $R\times S$ by the bimodule $M$ (Theorem~\ref{thm:triangular:jordan-super})
\cite[p.~3]{CheraghpourJafari2024}.\begin{theorem}[{\cite[Theorem~2.3, pp.~4--6]{CheraghpourJafari2024}}]\label{thm:triangular:jordan-super}
Let $d_0$ and $d_1$ be Jordan superderivations of degree $0$ and $1$ of the
upper triangular matrix ring $T=\begin{pmatrix}R&M\\0&S\end{pmatrix}$, with $R$
and $S$ $2$-torsion free. Then there are Jordan derivations $\delta_1$ of $R$
and $\delta_2$ of $S$, an additive map $g$ on $M$ with
$g(rm+ms)=rg(m)+\delta_1(r)m+g(m)s+m\delta_2(s)$, and an element $m^{*}\in M$
such that
\[
d_0\begin{pmatrix}r&m\\0&s\end{pmatrix}
=\begin{pmatrix}\delta_1(r)&g(m)\\0&\delta_2(s)\end{pmatrix},
\qquad
d_1\begin{pmatrix}r&m\\0&s\end{pmatrix}
=\begin{pmatrix}0&rm^{*}-m^{*}s\\0&0\end{pmatrix}.
\]
\end{theorem}

So the even part is a triangular Jordan derivation. The odd part is the inner
superderivation of one element \cite[p.~6]{CheraghpourJafari2024}. The corollary
is that every Jordan superderivation of the ring is a Jordan derivation
\cite[p.~6]{CheraghpourJafari2024}.

Two results come from other chapters and are not repeated here. The
$f$-derivation theorem of Benkovič concerns triangular algebras. It says that an
$f$-derivation is a derivation when the coefficients of $f$ have a nonzero sum
that is regular \cite[Theorem~1.1]{Benkovic2014}, and it is stated as
Theorem~\ref{thm:f:triangular}. The triangular-ring form is a unital ring with a
nontrivial idempotent $e$ such that $f=1-e$ and $fRe=0$. It appears in
Section~\ref{sec:higher-multi:tri-generalized} \cite[p.~1953]{Wang2016Bideriv}.

Recall that a triangular algebra has two square blocks of possibly different sizes. The
generalized matrix algebra of
Section~\ref{sec:matrix-triangular:generalized-matrix} keeps both off-diagonal
blocks, and the two bilinear maps that multiply them are what remains of the matrix
product.

The block upper triangular case has the same shape. Ghahramani, Ghosseiri and
Heidarizadeh prove that a linear map on $T(n_1,\dots,n_k)$ satisfying the Jordan
zero-product condition has the form $D+\alpha+\cdot\Delta(I)$, with $D$ a
derivation \cite[Abstract]{GhahramaniGhosseiriHeidarizadeh2020}.

Every multiplicative Lie triple derivation of a triangular 3-matrix ring has the standard form~\cite{WlMultiplicativeLieTriple2021}.

Every multiplicative Lie triple derivation of a triangular 3-matrix ring has the standard form~\cite{WlMultiplicativeLieTriple2021}.

A generalized Jordan derivation associated with a Hochschild 2-cocycle on a triangular algebra is shown to be a generalized derivation under stated conditions~\cite{WlGeneralizedJordanDerivationsAs2019}.

For a triangular algebra $\mathrm{Tri}(A,M,B)$ over a commutative ring with unity, Ashraf and Jabeen show that under stated assumptions every nonlinear generalized Lie triple higher derivation is of standard form: each component is a sum of an additive generalized higher derivation and a nonlinear functional vanishing on all Lie triple products~\cite{WlNonlinearLieTripleHigherDeriva2018}.

For a triangular algebra $\mathcal{A}=\mathrm{Tri}(A,M,B)$ over a commutative ring with unity, Ashraf, Jabeen and Parveen prove that every Jordan triple $(\sigma,\tau)$-higher derivation of $\mathcal{A}$ is a $(\sigma,\tau)$-higher derivation~\cite{WlOnJordanTripleHigherDerivation2018}.

Ashraf, Akhtar and Ansari describe the Lie and Jordan triple derivations and the generalized Lie and Jordan triple derivations of an arbitrary triangular algebra through a larger triangular algebra built from the maximal left and right rings of quotients, in which the given algebra sits with the same unity~\cite{WlGeneralizedLieJordanTripleDeri2021}.

Driss and L'Moufadal prove for a generalized triangular matrix algebra that every Jordan automorphism is an automorphism or an antiautomorphism and every Jordan derivation is a derivation~\cite{WlJordanAutomorphismsJordanDeriv2005}.

Ghahramani, Ghosseiri and Rezaei characterise the Jordan derivable maps of a 2-torsion free triangular ring by their local action on the pairs with $AB=0$ and on the standard idempotent~\cite{WlCharacterizingJordanDerivableM2022}.
